% ---
% RESUMOS
% ---

% RESUMO em português
\setlength{\absparsep}{18pt} % ajusta o espaçamento dos parágrafos do resumo
\begin{resumo}
  O rastreio contínuo de satélites em órbita baixa (LEO) exige sistemas de apontamento dinâmicos e precisos para mitigar perdas de enlace decorrentes da alta velocidade relativa e da diretividade das antenas.
  Este trabalho apresenta o desenvolvimento, modelagem cinemática, implementação e controle de um sistema de pedestal bi-axial (azimute e elevação) automatizado para o rastreio de satélites e recepção de dados meteorológicos de alta resolução (HRPT) na faixa de 1,7 GHz. A cinemática direta do mecanismo foi formulada via parâmetros de Denavit-Hartenberg (D-H). 
  Para garantir a estabilidade e suavidade da trajetória em tempo real com baixo custo computacional, implementou-se um filtro preditivo $\alpha-\beta-\gamma$ embarcado, cujos parâmetros de suavização foram otimizados via varredura sistemática e validados pelo critério de estabilidade de Jury. A arquitetura de hardware e controle foi estruturada em um microcontrolador dual-core ESP32 sob o sistema operacional FreeRTOS, segregando as tarefas de comunicação TCP/interface LCD e a malha de controle do filtro com acionamento dos motores de passo NEMA 17 via drivers DRV8825. A recepção e filtragem do sinal de radiofrequência foram realizadas por meio de uma antena helicone customizada, aliada a um módulo LNA/filtro SAW e um dispositivo SDR. Os resultados demonstram a funcionalidade mecânica e eletrônica do protótipo de baixo custo, garantindo cobertura hemisférica completa e suavização satisfatória na estimativa de posições orbitais.

  \vspace{\onelineskip}
  \textbf{Palavras-chave}: Rastreio de satélites. Pedestal bi-axial. Cinemática D-H. Filtro $\alpha-\beta-\gamma$. Sistema embarcado. ESP32.
\end{resumo}

% ABSTRACT in english
\begin{resumo}[Abstract]
 \begin{otherlanguage*}{english}
  Continuous tracking of Low-Earth Orbit (LEO) satellites requires dynamic and precise pointing systems to mitigate link losses resulting from high relative velocity and antenna directivity. 
  This work presents the development, kinematic modeling, implementation, and control of an automated bi-axial (azimuth and elevation) pedestal system for satellite tracking and high-resolution meteorological data reception (HRPT) at 1.7 GHz. The direct kinematics of the mechanism was formulated using Denavit-Hartenberg (D-H) parameters. 
  To ensure trajectory stability and smoothness in real-time with low computational overhead, an embedded predictive $\alpha-\beta-\gamma$ filter was implemented, whose smoothing parameters were optimized via grid search and validated using Jury's stability criterion. 
  The hardware and control architecture was structured on a dual-core ESP32 microcontroller running the FreeRTOS real-time operating system, segregating TCP communication and LCD interface tasks from the predictive filter loop and stepper motor driving (NEMA 17 via DRV8825 drivers). 
  Radio frequency signal reception and filtering were accomplished using a custom helicone antenna combined with an LNA/SAW filter module and a Software Defined Radio (SDR). The results demonstrate the mechanical and electronic feasibility of the low-cost prototype, ensuring complete hemispheric coverage and satisfactory smoothing of orbital position estimates.

  \vspace{\onelineskip}
  \textbf{Keywords}: Satellite tracking. Bi-axial pedestal. D-H kinematics. $\alpha-\beta-\gamma$ filter. Embedded system. ESP32.
 \end{otherlanguage*}
\end{resumo}